% Modernised reading — fonds Alexandre Grothendieck, shelfmark 123, pages 1-47.
%
% This is an INTERPRETATION, not a transcription. Everything the critical
% apparatus carried moves into footnotes. In any dispute of substance the
% transcription governs: whoever cites Grothendieck must cite
% batch-01.fr.tex, batch-02.fr.tex or batch-03.fr.tex.
%
% Pass: Opus 5.5 (claude-opus-5-5), 2026-10-03 — first pass, unchecked against
% the pages by a human. The folder was read whole: three batches, pages 1-47,
% all transcribed. Covers (3, 7, 14, 23, 36, 42) and blanks (20, 24, 32) carry
% no \page in the transcription; their titles head the sections here.
%
% Notation fixed for the whole document.
% — G is a reductive (mostly semisimple adjoint) group scheme over a base S;
%   g = Lie(G); T a maximal torus, t = Lie(T), W the Weyl group; h the Coxeter
%   number, z the order of the centre of the simply connected cover, written
%   frak{z} as a group; phi Euler's function.
% — Bor is the scheme of Borel subgroups everywhere. Pages 8-13 call it D; the
%   reading renames it, because from page 21 on D is a cyclic line, and the
%   folder never uses D for Borels again.
% — The Lie-algebra Grothendieck-Springer space is tilde-g (pages 4-6 write
%   underlined L), the group one tilde-G (pages 8-13 write round B, the
%   "universal Borel", which is the same scheme). J = t/W and I = T/W are kept.
% — M is the character lattice of T on pages 8-13 and a cocharacter lattice
%   on page 17. The reading writes X^*(T) and X_*(T) and footnotes the clash.
% — Two "principal" cocharacters: the Dynkin cocharacter of the principal
%   SL(2) (pages 15-19, 40), which is 2 rho-check, and the "hom. principal"
%   of page 43, which is rho-check, defined because G is adjoint. Fixed in the
%   spine section; the page-40 statement "all components 1" is corrected to 2.
%
% Substantive departures from the pages, each footnoted where it occurs:
% — p. 1: w c w^{-1} = c^h is corrected to c^n.
% — pp. 4-5: free action of W on t^rss is kept as the hypothesis it is
%   (question 0 of page 6 holds it open); the existence of f: g -> J is
%   stated as equivalent to Chevalley restriction, true in char. 0.
% — p. 9: "f normal" is stated under the hypothesis that the derived group is
%   simply connected, and G' normal under S normal.
% — p. 10, (iv): Pic(Bor) -> Pic(tilde-G) is an isomorphism only when
%   H^1(Gamma, X^*(T)) = 0 (as page 12 itself says); the ample cone is the
%   interior of P^+, not P^+.
% — p. 11: "T = prod_{k'/k} T'" is read as "T splits over k'".
% — p. 13, d): the finiteness of the cokernel of Pic(S) -> Pic(G^reg) needs S
%   regular, not only the injectivity the page proves.
% — p. 15, b): ad(e): g(0) -> g(-2) is corrected to g(0) -> g(2).
% — p. 18: the inference that a rational Dynkin-Kostant parabolic forces G
%   quasi-split is established only in the principal case.
% — p. 19: "g(1) = 0 iff principal" (blank filled as the transcription
%   suggests) is false; g(1) = 0 characterises even classes. Stated correctly.
% — p. 40, b): principal = all weights 2, not 1.
% — p. 44: "racines positives" corrected to "racines"; p. 45-46: "sauf une"
%   to "sauf au plus une".
%
% What the folder announces and never establishes: the answers to questions
% 0)-g) of page 6; the extension asked for in the margin of page 1; the count
% card(W)^2/(h z phi(h)) of the letter and of page 34 (asserted, not proved);
% the "à prouver" of page 30; the normaliser of a Kostant homomorphism
% (page 46); the paragraph "e)" of page 41, cut off after one sentence; the
% struck "e)" at the foot of page 13. The typed versos (27, 29, 31, 35, 39)
% belong to two other texts and are only named, as the transcription names
% them (RIGHTS.md).
\documentclass[11pt,a4paper]{article}
\input{../preamble/grothendieck}

\folder{123}
\pages{1}{47}
\dating{1969}
\watermark{Édition de démonstration\\interprétation personnelle de l'œuvre}
\foldertitle{[Autour de Kostant] : tapuscrits (s.d.), notes manuscrites (s.d.), lettre (1969) — lecture modernisée du dossier entier}

\begin{document}

\section*{Résumé}

\begin{resume}
Ce dossier est un relevé de lecture : Grothendieck relit, en 1969, deux
articles de Bertram Kostant sur les groupes de Lie, et les récrit dans sa propre
langue — celle des schémas, où tout énoncé doit valoir non seulement sur les
nombres complexes, mais au-dessus d'une base quelconque, en famille.

Le point de départ est une application très concrète. À une matrice carrée,
associons ses valeurs propres, prises comme un ensemble non ordonné —
c'est-à-dire son polynôme caractéristique. Les matrices qui ont mêmes valeurs
propres forment une « fibre » de cette application. Une fibre générique, où les
valeurs propres sont distinctes, est lisse : c'est une seule classe de matrices
diagonalisables. Mais la fibre où toutes les valeurs propres sont nulles —
les matrices nilpotentes — est singulière, et d'autant plus singulière que la
matrice est dégénérée. La même application existe pour tout groupe de Lie
semi-simple : c'est l'application de Kostant $\mathfrak{g}\to\mathfrak{t}/W$
pour l'algèbre de Lie, et celle de Steinberg $G\to T/W$ pour le groupe.

Le premier problème du dossier est de « lisser » cette famille de fibres. Le
cas le plus petit dit tout. Pour les matrices $2\times 2$ de trace nulle, la
fibre au-dessus de $s$ est la quadrique $a^{2}+bc=s$ : un hyperboloïde pour
$s\neq 0$, un cône pour $s=0$. On ne peut pas désingulariser le cône sans
abîmer ses voisins ; mais si l'on pose $s=t^{2}$ — si l'on choisit une racine
carrée, c'est-à-dire si l'on \emph{ordonne} les deux valeurs propres $\pm t$ —
la famille devient résoluble d'un seul coup. Dans le cas général, ordonner les
valeurs propres revient à passer de $\mathfrak{t}/W$ à $\mathfrak{t}$, et la
résolution s'obtient en adjoignant à chaque élément un sous-groupe de Borel qui
le contient — un « drapeau ». Les feuillets 4 à 13 montrent que cette
construction résout la famille changée de base, et prouvent, par un calcul de
groupes de Picard, qu'aucune résolution analogue n'existe \emph{sans} ce
changement de base. C'est ce que la littérature appelle la résolution simultanée
de Grothendieck.

Le second fil est arithmétique et combinatoire. Kostant avait isolé, dans tout
groupe simple, des objets privilégiés : le sous-groupe « principal » de type
$\mathrm{SL}(2)$, l'élément de Coxeter du groupe de Weyl — qui agit comme une
rotation d'ordre maximal $h$ —, et les éléments « cycliques », sommes des
vecteurs de racines simples et du vecteur de la plus basse racine. Grothendieck
les reprend un par un, les range en tableaux d'espaces homogènes avec leurs
stabilisateurs, et se demande ce qui en reste sur un corps qui n'est pas
algébriquement clos : par exemple, l'existence d'un sous-groupe principal défini
sur le corps équivaut à ce que le groupe soit quasi-déployé. Il interprète aussi
la théorie de Kostant en termes de filtrations du foncteur qui oublie la
structure d'une représentation — une première esquisse de ce qui deviendra la
théorie tannakienne des filtrations.

Au milieu du dossier, une lettre à Kostant, datée de Massy le 22 octobre 1969,
pose ses questions en clair : comment varie le rang de l'application de Kostant
sur les orbites nilpotentes, quelle est la dimension de la variété des Borels
contenant un élément donné, les singularités des fibres sont-elles
« rationnelles » ? Elle dit aussi d'où vient la relecture : d'un travail de
Brieskorn sur les singularités de Klein. La lettre donne la mesure de
l'attitude : relire une œuvre déjà classique parce qu'elle « appelle une
reconsidération soigneuse dans le cadre de la géométrie algébrique », puis poser
les questions qu'elle suggère, quitte à découvrir qu'elles sont connues —
« pour certaines, il serait scandaleux que les experts ne connaissent pas la
réponse ». La plupart de ces questions ont reçu depuis des réponses affirmatives.

Pour aller plus loin : résolution de Grothendieck--Springer, application de
Steinberg, section de Kostant, théorème de Jacobson--Morozov, éléments réguliers
de Springer, orbite sous-régulière et théorème de Brieskorn--Slodowy.

\keywords{Grothendieck-Springer resolution, simultaneous resolution, adjoint quotient, Steinberg map, Chevalley restriction theorem, Kostant section, regular element, Picard group of flag variety, Coxeter element, Coxeter number, regular elements of reflection groups, principal three-dimensional subgroup, sl2-triple, Jacobson-Morozov theorem, weighted Dynkin diagram, Jacobson-Morozov parabolic, quasi-split group, Tannakian filtration, Springer fibre, subregular orbit, Klein singularity, rational singularities, cyclic element, extended affine Weyl group}
\end{resume}

\pagerange{1}{47}
\section*{Le fil du dossier, et ses conventions}

Quarante-sept feuillets, dont six chemises qui ne portent qu'un titre de sa main
(feuillets 3, 7, 14, 23, 36, 42) et trois blancs (20, 24, 32). Trois sortes de
papier s'y mêlent : deux tapuscrits de sa main, en français, corrigés à
l'encre (feuillets 1--2 et 15--19) ; une lettre dactylographiée en anglais
(21--22) ; des notes manuscrites, dont plusieurs sur le verso de feuillets
dactylographiés étrangers au dossier.\footnote{Les versos 27, 29, 31, 35 et 39
appartiennent à deux autres textes : un exposé sur les topos (§~3.1, exactitude
des foncteurs $u^{*}$ et $u_{*}$ d'un morphisme de topos) et un paragraphe
« Accouplements définis par une biextension ». Ils ne touchent pas aux notes sur
Kostant ; la transcription les résume et en relève les ajouts manuscrits,
l'édition ne fait que les nommer.} La date de l'inventaire, 1969, est celle de
la lettre ; le reste est sans date, mais tout se tient autour d'elle, et
plusieurs feuillets portent des reprises à l'encre brune, plus tardives.

Le dossier se lit en huit stations.

\begin{itemize}
\item \textbf{La transformation de Coxeter} (feuillets 1--2) : ce que Kostant
  dit des valeurs propres d'un élément de Coxeter, et les « droites
  cycliques » du groupe de Weyl.
\item \textbf{Lissification des morphismes de Kostant} (3--6) : la résolution
  simultanée pour l'algèbre de Lie, et une liste de questions.
\item \textbf{Lissification des morphismes de Steinberg} (7--13) : la même
  construction pour le groupe, avec un théorème, un lemme sur les groupes de
  Picard, et la preuve qu'on ne peut résoudre sans changer de base.
\item \textbf{Résumé des théorèmes de conjugaison de Kostant} (14--19) : les
  sous-groupes de type $\mathrm{SL}(2)$, les éléments de Dynkin, et leur
  interprétation par les filtrations du foncteur fibre.
\item \textbf{La lettre à Kostant} (21--22).
\item \textbf{Diagrammes d'espaces homogènes} (23--35) : éléments de Kostant
  et de Coxeter, tores en apposition, droites cycliques, rigidifications.
\item \textbf{Liste d'espaces homogènes remarquables} (36--41), suivie de
  « Résultats de Kostant ».
\item \textbf{Résumé de quelques résultats de Kostant} (42--47), mis au
  propre : homomorphismes principaux, homomorphismes de Kostant
  $\mu_{h}\to G$, graduation de $\mathfrak{g}$, droites cycliques.
\end{itemize}

Les deux premières stations « géométriques » (3--13) et les autres
« combinatoires » ne se citent presque pas l'une l'autre ; ce qui les lie est
Kostant, dont les deux articles sont nommés dans la lettre : \emph{The principal
three-dimensional subgroup and the Betti numbers of a complex simple Lie group}
(Amer. J. Math., 1959) et \emph{Lie group representations on polynomial rings}
(Amer. J. Math., 1963). Le dossier contient deux résumés de Kostant, le
tapuscrit des feuillets 15--19 et le manuscrit des feuillets 43--47 : ce sont
deux moutures distinctes, sur des sujets voisins mais non identiques, et
l'édition les garde séparées.

\textbf{Conventions, valables pour tout le dossier.} $G$ est un schéma en
groupes réductif sur une base $S$, le plus souvent semi-simple adjoint ;
$\mathfrak{g}=\operatorname{Lie}(G)$ ; $T$ est un tore maximal,
$\mathfrak{t}=\operatorname{Lie}(T)$, $W$ le groupe de Weyl ; $r$ est le rang,
$N$ le nombre de racines positives, $n=\dim G=2N+r$. On note $h$ le nombre de
Coxeter, $\psi$ la plus grande racine, $\Sigma$ l'ensemble des racines simples,
$\mathfrak{z}$ un groupe isomorphe au centre du revêtement simplement connexe
$\widetilde{G}$ et $z$ son ordre, $\varphi$ l'indicatrice d'Euler. On écrit
$X^{*}(T)$ et $X_{*}(T)$ pour les réseaux des caractères et des
cocaractères.\footnote{Le dossier écrit $M$ pour l'un \emph{et} pour l'autre :
$M=\widehat{T}$, réseau des caractères, aux feuillets 8 à 13 ;
$M=\operatorname{Hom}(\mathbf{G}_{m},T_{0})$, réseau des cocaractères, au
feuillet 17. L'édition lève l'ambiguïté.}

Trois homonymies sont à fixer dès maintenant.

\begin{itemize}
\item $\mathrm{Bor}$ désigne partout le schéma des sous-groupes de Borel. Les
  feuillets 8 à 13 l'appellent $D$ ; à partir du feuillet 21, $D$ est une
  \emph{droite cyclique}. L'édition réserve $D$ aux droites.
\item Le « Borel universel » des feuillets 4 et 8, vu comme schéma des couples
  $(g,B)$ avec $g\in B$, est noté $\widetilde{G}$ ; son analogue pour l'algèbre
  de Lie, le schéma des couples $(x,\mathfrak{b})$ avec $x\in\mathfrak{b}$, est
  noté $\widetilde{\mathfrak{g}}$.\footnote{Les feuillets écrivent
  $\mathfrak{B}$ (un $B$ rond) et $\underline{\mathfrak{L}}$ (un $L$ rond
  souligné). Les notations $\widetilde{\mathfrak{g}}$ et $\widetilde{G}$ sont
  les notations actuelles ; à ne pas confondre avec le revêtement simplement
  connexe, que l'édition ne nomme $\widetilde{G}$ que dans la lettre, comme
  elle.}
\item \emph{Principal} a deux sens. Au feuillet 15 et au feuillet 40, le
  cocaractère attaché à un sous-groupe principal $\mathrm{SL}(2)\to G$ est le
  cocaractère de Dynkin, qui vaut $2$ sur chaque racine simple : c'est
  $2\rho^{\vee}$. Au feuillet 43, un « homomorphisme principal »
  $\mathbf{G}_{m}\to G$ vaut $1$ sur chaque racine simple : c'est
  $\rho^{\vee}$, somme des copoids fondamentaux, qui n'existe comme
  cocaractère que parce que $G$ est adjoint. Le second est la « moitié » du
  premier.
\end{itemize}

\pagerange{1}{2}
\section*{Compléments sur la transformation de Coxeter (feuillets 1 et 2)}

Deux feuillets dactylographiés, intitulés « Compléments sur transformation de
Coxeter d'un système de racines irréductible (Kostant) », avec en marge un
renvoi à l'article de 1959 et à son théorème 9.2.\footnote{« irréductible »
est ajouté à la main dans le titre. Le tapuscrit est de lui : il est en
français, cite Kostant à la troisième personne et est corrigé de sa main, qui
complète aussi les lettres grecques que la machine n'avait pas.} Soit $R$ un
système de racines irréductible dans un espace réel $V$, $W$ son groupe de
Weyl, $m_{1}\leqslant\dots\leqslant m_{r}$ ses exposants et $h=m_{r}+1$ son
nombre de Coxeter. On fait agir $W$ sur $V\otimes\mathbf{C}$.

\textbf{a)} Toute valeur propre d'un $w\in W$ est une racine de l'unité dont
l'ordre divise l'un des degrés $m_{i}+1$, donc est au plus $h$. Une valeur
propre d'ordre exactement $h$ existe si et seulement si $w$ est un élément de
Coxeter ; le sous-espace propre correspondant est alors une droite, et ses
vecteurs non nuls sont réguliers (n'appartiennent à aucun hyperplan de
réflexion).\footnote{La parenthèse « (qui est de dim.\ 1) » est une note
marginale de lecture incertaine. Le tapuscrit attribue l'ensemble à Kostant.
La forme générale de la première phrase — un $\zeta$ primitif d'ordre $d$ est
valeur propre d'un élément de $W$ si et seulement si $d$ divise un degré — est
aujourd'hui citée d'après Springer, \emph{Regular elements of finite reflection
groups} (Invent. Math., 1974), postérieur au dossier ; l'édition ne tranche pas
ce que Kostant en énonçait en 1959.}

\textbf{Corollaire 1.} Soit $c$ un élément de Coxeter et $\mu=\langle
c\rangle\simeq\mathbf{Z}/h\mathbf{Z}$. Alors $\mu$ est son propre
centralisateur dans $W$, et si $\nu$ est le normalisateur de $\mu$,
l'homomorphisme $\nu/\mu\to\operatorname{Aut}(\mu)\simeq
(\mathbf{Z}/h\mathbf{Z})^{*}$ est un isomorphisme. Autrement dit, pour tout
$n$ premier à $h$, il existe $w\in W$, unique modulo $\mu$, tel que
$wcw^{-1}=c^{n}$.\footnote{Le tapuscrit porte $wcw^{-1}=c^{h}$, qui serait
$=1$ ; le sens, et la phrase de la démonstration (« si $n$ est premier à $h$,
alors $c^{n}$ est conjugué de $c$ »), demandent $c^{n}$, que l'édition écrit.
Deux notes marginales accompagnent l'énoncé : « Quelle extension de
$(\mathbf{Z}/h\mathbf{Z})^{*}$ par $\mathbf{Z}/h\mathbf{Z}$ obtient-on ? » —
question que le dossier ne résout pas, et que la lettre reprend sous une autre
forme — et, d'une lecture incertaine, « échanger avec cor.\ 3 ? », alors
qu'aucun corollaire n'est numéroté 3.}

\emph{Démonstration.} Soit $D$ la droite propre de $c$ pour une racine
primitive $h$-ième de l'unité.\footnote{Le tapuscrit justifie $\dim D=1$ par une
« prop.\ 30 » dont la source n'est pas indiquée.} Le centralisateur
$C_{W}(c)$ normalise $D$, d'où un homomorphisme
$C_{W}(c)\to\operatorname{Aut}(D)=\mathbf{C}^{*}$ ; son image est
un groupe cyclique fini qui contient l'image, isomorphe, de $\mu$, donc d'ordre
multiple de $h$ ; comme ses éléments sont des valeurs propres d'éléments de
$W$, leur ordre est au plus $h$, et l'image est d'ordre $h$. On obtient une
rétraction $C_{W}(c)\to\mu$, dont le noyau fixe $D$ point par point
; $D$ étant régulière et $W$ agissant simplement transitivement sur les
chambres, ce noyau est trivial. Reste que $c^{n}$ est conjugué de $c$ pour $n$
premier à $h$, que le tapuscrit tire encore de Kostant.

\textbf{Corollaire.} Soit $x$ un vecteur propre de $c$, de valeur propre
$\lambda$. Alors $x$ est régulier si et seulement si $\lambda$ est une racine
primitive $h$-ième de l'unité. La suffisance est a). Pour la nécessité, si
$h'$ est l'ordre de $\lambda$, $c^{h'}$ fixe $x$, donc $c^{h'}=1$ puisque $x$
est régulier, et $h'=h$.

\textbf{Corollaire.} Soit $D\subset V\otimes\mathbf{C}$ une droite régulière
normalisée par un élément de Coxeter $c$.\footnote{La frappe superpose deux
débuts, « Soit D une » et « Considérons », tous deux recouverts ; l'édition lit
« Soit $D$ une droite régulière ».} Alors le normalisateur de $D$ dans $W$ est
$\langle c\rangle$, et pour chaque racine primitive $h$-ième $\zeta$ il existe
un unique élément de Coxeter normalisant $D$ et y agissant par $\zeta$. Ces
droites — les \emph{droites cycliques} de Kostant — forment une seule orbite
sous $W$, et leur nombre est $\operatorname{card}(W)/h$. Le groupe
$(\mathbf{Z}/h\mathbf{Z})^{*}$ agit librement sur l'ensemble des droites
cycliques, en commutant à $W$ ; deux droites cycliques sont dans la même orbite
de ce groupe si et seulement si elles ont même stabilisateur dans $W$.

Ce sont, au niveau du groupe de Weyl, les objets dont les feuillets 21 à 47
construiront les relèvements dans $G$ et dans $\mathfrak{g}$.

\pagerange{3}{6}
\section*{Lissification des morphismes de Kostant (feuillets 3 à 6)}

\pagerange{4}{5}
\subsection*{Les données, et le morphisme de Kostant}

Le feuillet 4 est une liste de notations. $G$ est un schéma en groupes
semi-simple \emph{adjoint} sur $S$, $\mathfrak{g}$ son algèbre de Lie, vue
comme schéma vectoriel sur $S$ ; $T$ est « le » tore maximal abstrait, défini
par descente à partir d'un couple de Killing, $\mathfrak{t}$ son algèbre de
Lie, $W$ « le » groupe de Weyl. On pose
\[
  \mathfrak{J}=\mathfrak{t}/W,\qquad \mathfrak{J}'=\mathfrak{t},
\]
de sorte que $\mathfrak{J}'\to\mathfrak{J}$ est un revêtement ramifié. On note
$\mathrm{Tor}$ le schéma des tores maximaux, $\mathrm{Kill}$ celui des couples
de Killing (tore maximal, Borel le contenant), $\mathrm{Bor}$ celui des Borels,
et $\widetilde{\mathfrak{g}}\to\mathrm{Bor}$ le fibré vectoriel des couples
$(x,\mathfrak{b})$, $x\in\mathfrak{b}$. Il porte deux morphismes : la
projection $q\colon\widetilde{\mathfrak{g}}\to\mathfrak{g}$, propre, et
$\tilde f'\colon\widetilde{\mathfrak{g}}\to\mathfrak{t}$, qui envoie
$(x,\mathfrak{b})$ sur l'image de $x$ dans
$\mathfrak{b}/[\mathfrak{b},\mathfrak{b}]\simeq\mathfrak{t}$, et qui est
lisse.

Les exposants $\mathrm{rss}$ désignent les ouverts d'éléments réguliers
semi-simples. On a un isomorphisme canonique
$\widetilde{\mathfrak{g}}^{\mathrm{rss}}\simeq
\mathrm{Kill}\times_{S}\mathfrak{t}^{\mathrm{rss}}$ — un élément régulier
semi-simple détermine son tore centralisateur, et le Borel choisi en fait un
couple de Killing — et un carré cartésien

\begin{tikzcd}
\mathfrak{g}^{\mathrm{rss}} \arrow[d] & \widetilde{\mathfrak{g}}^{\mathrm{rss}} \arrow[l] \arrow[d] \\
\mathrm{Tor} & \mathrm{Kill} \arrow[l]
\end{tikzcd}

où la flèche de gauche associe à $x$ son centralisateur. Le groupe $W$ opère
sur $\mathrm{Kill}$, qui est un torseur sous $W$ au-dessus de $\mathrm{Tor}$,
d'où une opération sur $\widetilde{\mathfrak{g}}^{\mathrm{rss}}$, diagonale
sur $\mathrm{Kill}\times\mathfrak{t}^{\mathrm{rss}}$, et
$\mathfrak{g}^{\mathrm{rss}}=\widetilde{\mathfrak{g}}^{\mathrm{rss}}/W$. On en
tire un morphisme naturel
$f^{\mathrm{rss}}\colon\mathfrak{g}^{\mathrm{rss}}\to
\mathfrak{J}^{\mathrm{rss}}=\mathfrak{t}^{\mathrm{rss}}/W$.

\textbf{Hypothèse (L).} $W$ opère librement sur
$\mathfrak{t}^{\mathrm{rss}}$, ce qui revient à dire qu'il opère fidèlement sur
les fibres géométriques de $\mathfrak{t}$. Sous (L), le carré formé de
$f^{\mathrm{rss}}$, de $q^{\mathrm{rss}}$, de
$\tilde f'^{\mathrm{rss}}$ et de $\mathfrak{t}^{\mathrm{rss}}\to
\mathfrak{J}^{\mathrm{rss}}$ est cartésien.\footnote{Le feuillet écrit d'abord
« trivialement », biffé, puis « librement », et ajoute que c'est le cas quand
les caractéristiques résiduelles sont toutes $\neq 2$ — avec un « si » de
lecture douteuse. L'édition garde (L) comme une hypothèse, ce que le dossier
fait lui-même : c'est la question 0) du feuillet 6. Elle est vraie en
caractéristique nulle ; elle tombe en caractéristique $2$ déjà pour
$\mathrm{PGL}_{2}$, où $W=\{\pm 1\}$ opère trivialement sur $\mathfrak{t}$.}

Le morphisme $f^{\mathrm{rss}}$ se prolonge (de façon unique) en un morphisme
\[
  f\colon\mathfrak{g}\longrightarrow\mathfrak{J}=\mathfrak{t}/W,
\]
le \emph{morphisme de Kostant} — aujourd'hui le quotient adjoint
$\mathfrak{g}\to\mathfrak{g}/\!\!/G$. Une note marginale le dit : l'anneau de
$\mathfrak{J}$ est l'anneau des invariants de celui de $\mathfrak{g}$. Le
prolongement revient, sur un corps, à la surjectivité de la restriction de
Chevalley $k[\mathfrak{g}]^{G}\to k[\mathfrak{t}]^{W}$ ; elle est classique en
caractéristique nulle.\footnote{La démonstration du feuillet 5 est une
esquisse : on se ramène à $S=\operatorname{Spec}\mathbf{Z}$ et $G$ déployé, et
comme $\mathfrak{g}$ est normal et $\mathfrak{J}$ affine, « il suffit de le voir
sur les fibres géométriques ; là c'est connu (Kostant…) ». Kostant travaille
sur $\mathbf{C}$ ; en petite caractéristique, la restriction de Chevalley peut
cesser d'être surjective, et l'édition ne prend pas à son compte le
prolongement sur $\mathbf{Z}$. La note marginale du feuillet 5 est en partie
effacée ; seul son début se lit sûrement.}

\pagerange{5}{6}
\subsection*{La résolution simultanée}

On pose $\mathfrak{g}'=\mathfrak{g}\times_{\mathfrak{J}}\mathfrak{J}'$ :
c'est la famille des fibres de $f$ rendue par le changement de base qui ordonne
les « valeurs propres ». On obtient le diagramme

\begin{tikzcd}
\mathfrak{g} \arrow[d, "f"'] & \mathfrak{g}' \arrow[l] \arrow[d, "f'"] & \widetilde{\mathfrak{g}} \arrow[l, "p"'] \arrow[dl, "\tilde f'"] \arrow[r] & \mathrm{Bor} \\
\mathfrak{J} & \mathfrak{J}' \arrow[l] & &
\end{tikzcd}

où $p=(q,\tilde f')$ est projectif et $\tilde f'$ lisse.\footnote{Le feuillet
trace aussi la flèche composée $q$ de $\widetilde{\mathfrak{g}}$ vers
$\mathfrak{g}$, en arc ; l'édition l'omet, puisqu'elle vaut le composé. Le but
de la dernière flèche est écrit à l'encre pâle et se lit « $D$ » : c'est le
schéma de base de $\widetilde{\mathfrak{g}}$, que l'édition nomme
$\mathrm{Bor}$.} Sous (L), $p$ est birationnel : il induit un isomorphisme
$\widetilde{\mathfrak{g}}^{\mathrm{rss}}\simeq\mathfrak{g}'^{\mathrm{rss}}$.
C'est la \emph{résolution simultanée} : $\tilde f'$ est une famille lisse au-dessus
de $\mathfrak{t}$, et elle résout fibre à fibre la famille $f'$.\footnote{C'est
la construction que la littérature appelle résolution simultanée de
Grothendieck, ou de Grothendieck--Springer. Elle a été rendue publique par
Brieskorn (congrès de Nice, 1970), qui l'attribue à Grothendieck, puis exposée
par Slodowy (\emph{Simple singularities and simple algebraic groups}, 1980).
Springer avait publié en 1969 la résolution du cône nilpotent, qui en est la
fibre en $0$.}

Par réduction à la base $\mathbf{Z}$ ou $\mathbf{Z}[1/2]$ et par le Main
Theorem de Zariski, le plus grand ouvert de $\mathfrak{g}'$ au-dessus duquel
$p$ est un isomorphisme est formé des $x'$ dont la fibre $p^{-1}(x')$ est finie :
c'est l'image inverse de l'ouvert $\mathfrak{g}^{\mathrm{reg}}$ des éléments
\emph{réguliers}, ceux qui ne sont contenus que dans un nombre fini de
sous-algèbres de Borel.\footnote{L'argument suppose $\mathfrak{g}'$ normal, ce
que le feuillet 6 ramène, dans la question c), à la normalité des fibres de
$f$ ; en caractéristique nulle, c'est Kostant (1963).} D'après Springer
(\emph{Publ. Math. IHÉS} 30, 1966, 5.3), cet ouvert rencontre la fibre
nilpotente : il existe des éléments nilpotents réguliers. Le feuillet prévoit
d'en déduire qu'il rencontre toutes les fibres de $f$, et ajoute plus tard,
souligné deux fois : « \textbf{Non} en car.\ 2, cas $A_{1}$, $B_{2}$… peut-être
d'autres ».

\textbf{Questions.} Le feuillet 6 se clôt sur une liste, que le dossier ne
résout pas.

\begin{enumerate}
\item[0)] $W$ opère-t-il fidèlement sur $\mathfrak{t}^{\mathrm{rss}}$ ?
\item[a)] $\mathfrak{J}$ est-il régulier, c'est-à-dire $\mathfrak{J}'$ plat
  sur $\mathfrak{J}$ ?
\item[b)] $f$ est-il plat ? En fait, $f$ est plat si et seulement si
  $\mathfrak{J}$ est régulier et $f$ équidimensionnel.\footnote{L'équivalence
  est juste : $\mathfrak{g}$ est régulier, donc un morphisme plat
  $\mathfrak{g}\to\mathfrak{J}$ force $\mathfrak{J}$ régulier, et
  réciproquement un morphisme équidimensionnel d'un schéma de Cohen--Macaulay
  vers un schéma régulier est plat (EGA~IV~6.1.5). En marge : « Oui pour les
  “bons” nombres premiers ? ».}
\item[c)] Les fibres géométriques de $f$ sont-elles intègres, adhérences
  d'une orbite, à complémentaire de codimension $\geqslant 2$, normales ? Si
  oui, $f'$ est normal, donc $\mathfrak{g}'$ l'est, et
  $\widetilde{\mathfrak{g}}\to\mathfrak{g}'\to\mathfrak{g}$ est la
  factorisation de Stein de $q$.
\item[d)] Les points de lissité de $f$ sont-ils les points réguliers ?
\item[e)] Définir une stratification orbitale de $\mathfrak{g}$ sur
  $\mathbf{Z}$, la comparer à celle du groupe, et aux orbites unipotentes ;
  questions analogues pour $\widetilde{G}$ et $\widetilde{\mathfrak{g}}$.
\item[f)] Reprendre la théorie de Kostant : l'ouvert des points réguliers
  a-t-il une section ?
\item[g)] Étudier les fibres de $p$, c'est-à-dire la variété des
  sous-algèbres de Borel contenant un élément donné : quelle est sa dimension,
  en relation avec celle de l'orbite ?
\end{enumerate}

Sur $\mathbf{C}$, Kostant (1963) avait déjà répondu à a)--d) et f) : $f$ est
plat, ses fibres sont des intersections complètes normales et irréductibles,
adhérences de l'orbite régulière qu'elles contiennent, les points lisses de $f$
sont les points réguliers, et l'ouvert régulier admet une section — la
\emph{section de Kostant}. Les questions portent donc sur une base
quelconque.\footnote{Le feuillet est très chargé ; c), d) et e) ne se lisent que
par fragments, et l'édition n'en donne que le sens assuré. f) et g) sont dans
la marge, à l'encre brune. La question g) est celle que la lettre à Kostant
reprend (feuillet 22).}

\pagerange{7}{13}
\section*{Lissification des morphismes de Steinberg (feuillets 7 à 13)}

\pagerange{8}{9}
\subsection*{Les données, et le morphisme de Steinberg}

$G$ est maintenant un schéma en groupes réductif sur $S$, $\mathrm{Bor}$ le
schéma de ses Borels, $\widetilde{G}\subset G\times_{S}\mathrm{Bor}$ le Borel
universel, de radical unipotent $\widetilde{G}_{u}$.\footnote{Le feuillet
appelle $D$ le schéma des Borels et $\mathfrak{B}$ le Borel universel ; voir
les conventions.} Le tore $\widetilde{G}/\widetilde{G}_{u}$ sur
$\mathrm{Bor}$ descend en un tore $T$ sur $S$, de réseau des caractères
$X^{*}(T)$, avec un groupe de Weyl $W$ sur $S$ qui opère sur l'un et l'autre ;
on pose $I=T/W$. Sur le schéma $K=\mathrm{Kill}(G)$ des couples de Killing, le
tore universel s'identifie à $T\times_{S}K$ et le groupe de Weyl universel à
$W\times_{S}K$.

Il existe un unique morphisme
\[
  f\colon G\longrightarrow I=T/W,
\]
invariant par automorphismes intérieurs, dont la restriction à un tore maximal
$T_{0}$ contenu dans un Borel $B_{0}$ est la projection
$T_{0}\to T_{0}/W_{0}\simeq I$ : c'est le \emph{morphisme de Steinberg}.
La marge pose la question de son existence sans hypothèse de simple connexité,
et répond « Oui » en renvoyant à Steinberg.\footnote{Sur un corps
algébriquement clos, la restriction $k[G]^{G}\to k[T]^{W}$ est en effet un
isomorphisme pour tout groupe réductif connexe (Steinberg, \emph{Regular
elements of semisimple algebraic groups}, Publ. Math. IHÉS 25, 1965). La
même note demande si $f$ est « normal en tous cas », et objecte aussitôt, entre
crochets : « Mais alors il faudrait que $I$ soit lisse$/S$ en tous cas ! »
L'objection est fondée : un morphisme plat de source régulière a un but
régulier. Quand le groupe dérivé est simplement connexe, $I$ est lisse et $f$
est plat à fibres normales (Steinberg) ; en général, non.}

On définit $G'$ par le carré cartésien de gauche ci-dessous, et un morphisme
$p\colon\widetilde{G}\to G'$ par ses deux composantes :
$h(g,B)=g$, et $\tilde f(g,B)=$ l'image de $g$ dans $B/B_{u}\simeq T$.

\begin{tikzcd}
G \arrow[d, "f"'] & G' \arrow[l] \arrow[d, "f'"] & \widetilde{G} \arrow[l, "p"'] \arrow[dl, "\tilde f"] \\
I & T \arrow[l, "\pi"] &
\end{tikzcd}

L'égalité $\pi\tilde f=fh$ se vérifie sur l'ouvert dense
$\widetilde{G}^{\mathrm{rs}}$ des $(g,B)$ avec $g$ régulier semi-simple, où
elle est évidente. Le morphisme $h$ est propre ($\mathrm{Bor}$ l'est sur $S$),
donc $p$ l'est aussi ; $\tilde f$ est lisse. Si de plus $f$ est normal — c'est
le cas quand le groupe dérivé de $G$ est simplement connexe — et si $S$ est
normal, alors $f'$ est normal et $G'$ est normal.\footnote{Le feuillet 9 dit
seulement : « Utilisant le fait que $f$ est normal, donc $f'$ normal, on trouve
que $G'$ est normal. » Les deux hypothèses entre tirets sont de l'édition ; la
première est celle sur laquelle la marge du feuillet 8 hésitait. Le feuillet
ajoute aussi que $p$ est propre « comme on verra qu'il est birationnel » ; la
propreté ne dépend pas de la birationalité.}

On note $G^{\mathrm{reg}}$ l'ouvert des éléments réguliers — c'est l'ouvert de
lissité de $f$ — et $G'^{\mathrm{reg}}$, $\widetilde{G}^{\mathrm{reg}}$ ses
images inverses.

\pagerange{10}{11}
\subsection*{Le théorème (feuillet 10) et sa démonstration}

\textbf{Théorème.} Sous les hypothèses précédentes ($S$ normal, $G'$ normal) :

\begin{enumerate}
\item[(i)] $p\colon\widetilde{G}\to G'$ est projectif et
  $\tilde f\colon\widetilde{G}\to T$ est lisse.
\item[(ii)] $p$ induit un isomorphisme
  $\widetilde{G}^{\mathrm{reg}}\xrightarrow{\sim}G'^{\mathrm{reg}}$, et
  $G'^{\mathrm{reg}}$ est le plus grand ouvert au-dessus duquel $p$ est un
  isomorphisme.
\item[(iii)] $\widetilde{G}-\widetilde{G}^{\mathrm{reg}}$ est de codimension
  $\geqslant 2$ fibre par fibre, et $G'-G'^{\mathrm{reg}}$ a toutes ses
  composantes de codimension $3$ fibre par fibre sur $S$.
\item[(iv)] Si $S$ est localement noethérien et normal, les restrictions
  donnent des isomorphismes
  \[
    \operatorname{Pic}(\widetilde{G})\xrightarrow{\sim}
    \operatorname{Pic}(\widetilde{G}^{\mathrm{reg}})\xleftarrow{\sim}
    \operatorname{Pic}(G'^{\mathrm{reg}}),
  \]
  l'homomorphisme $\operatorname{Pic}(\mathrm{Bor})\to
  \operatorname{Pic}(\widetilde{G})$ est injectif, et il est bijectif lorsque
  $H^{1}(\Gamma,X^{*}(T))=0$, où $\Gamma$ est l'image du groupe fondamental de
  $S$ (en un point générique) dans $\operatorname{Aut}X^{*}(T)$ — en
  particulier lorsque $T$ est déployé. En faisceautisant pour la topologie
  étale, $\underline{\mathrm{Pic}}_{\mathrm{Bor}/S}$ s'identifie au réseau $P$
  des poids, $X^{*}(T)\to P$ étant injectif si $G$ est semi-simple, et un
  faisceau inversible sur $\mathrm{Bor}$ est ample si et seulement si son image
  inverse sur $\widetilde{G}$ l'est ; le cône ample est formé des poids
  strictement dominants.
\end{enumerate}

L'énoncé du feuillet 10 donne
$\operatorname{Pic}(\mathrm{Bor})\simeq\operatorname{Pic}(\widetilde{G})$ sans
condition, avec en marge une condition rajoutée après coup ; le feuillet 12
montre que la bijectivité tombe quand $T$ est tordu. L'édition met la condition
dans l'énoncé.\footnote{La note marginale, en partie illisible, lit : « Si $T$
est trivial, ou génériquement si l'image de $\pi_{1}(S,\xi)$ dans
$\operatorname{Ext}(G_{\xi})$ est telle que $H^{1}(\Gamma,M)=0$ » ; « on a une
réciproque si $S$ régulier ». Le feuillet écrit aussi que le cône ample « est
le cône des poids $P^{+}$ » ; un fibré en droites sur la variété des drapeaux
est ample exactement pour les poids \emph{strictement} dominants, $P^{+}$ étant
le cône des fibrés engendrés par leurs sections. Le mot qui précède « pour
$\underline{\mathrm{Pic}}_{D/S}$ » est illisible.}

\emph{Démonstration.} (i) et (iii) sont connus ; pour (iii), Steinberg (le
lieu non régulier de $G$ est de codimension $3$, et son image inverse dans
$\widetilde{G}$, où les fibres au-dessus des éléments sous-réguliers sont des
courbes, est de codimension $2$).

(ii) On a d'abord $\widetilde{G}^{\mathrm{rs}}\xrightarrow{\sim}
G'^{\mathrm{rs}}$ au-dessus de l'ouvert régulier semi-simple. Par Steinberg,
les fibres de $\widetilde{G}^{\mathrm{reg}}\to G'^{\mathrm{reg}}$ sont finies ;
le morphisme est propre, birationnel, et $G'^{\mathrm{reg}}$ est normal : le
Main Theorem de Zariski conclut.\footnote{La maximalité de
$G'^{\mathrm{reg}}$ n'est pas démontrée sur la page ; elle tient à ce qu'un
élément non régulier est contenu dans une famille de dimension positive de
Borels. Une note à l'encre brune, en partie illisible, ajoute un renseignement
tiré de Steinberg sur les centralisateurs.}

(iv) $S$ étant normal, $\widetilde{G}$ l'est, et l'injectivité de
$\operatorname{Pic}(\widetilde{G})\to\operatorname{Pic}(\widetilde{G}^{\mathrm{reg}})$
tient à la profondeur, puisque le complémentaire est de codimension
$\geqslant 2$ fibre par fibre. La surjectivité de
$\operatorname{Pic}(\mathrm{Bor})\to\operatorname{Pic}(\widetilde{G}^{\mathrm{reg}})$
se déduit de EGA~IV~21.4.13, compte tenu du lemme suivant ; les deux ensemble
donnent les isomorphismes annoncés.\footnote{Le renvoi est celui de la page,
avec « (EGA IV 53) » entre parenthèses, que l'édition ne s'explique pas.}

\textbf{Lemme.} Soit $G$ semi-simple simplement connexe sur un corps $k$, $B$
un Borel, $T=B/B_{u}$. Alors
\[
  \operatorname{Pic}(B)=\operatorname{Pic}(B^{\mathrm{reg}})
  =\operatorname{Pic}(T)=H^{1}(\Gamma,X^{*}(T)),
\]
où $\Gamma$ est le groupe de Galois d'une extension galoisienne finie $k'$ qui
déploie $T$. En effet, $B\to T$ est un torseur sous le groupe unipotent
déployé $B_{u}$, d'où $\operatorname{Pic}(B)=\operatorname{Pic}(T)$ ; la suite
spectrale de Hochschild--Serre pour $T_{k'}$ déployé donne
\[
  0\to H^{1}\bigl(\Gamma,\ k'^{*}\times X^{*}(T)\bigr)\to\operatorname{Pic}(T)
  \to H^{0}\bigl(\Gamma,\operatorname{Pic}(T_{k'})\bigr)=0,
\]
puisque les inversibles de $T_{k'}$ sont $k'^{*}\times X^{*}(T)$ ; et
$H^{1}(\Gamma,k'^{*})=0$ par le théorème 90 de Hilbert.\footnote{Le feuillet
écrit « $T=\prod_{k'/k}T'$, $T'$ un tore sur une extension galoisienne finie
déployée » : tout tore n'est pas une restriction des scalaires, et la
démonstration n'utilise que le fait que $T$ se déploie sur $k'$, qu'on écrit.
L'égalité avec $\operatorname{Pic}(B^{\mathrm{reg}})$ vient de la codimension
$\geqslant 2$ du lieu non régulier dans $B$, qui est lisse. Plusieurs termes
du lemme sont ajoutés en couches au-dessus de la ligne.}

Si $T$ est tordu, ce groupe peut être non nul ; prenant pour $B$ la fibre
générique de $\widetilde{G}$ sur $\mathrm{Bor}$, on voit que
$\operatorname{Pic}(\mathrm{Bor})\to\operatorname{Pic}(\widetilde{G})$ n'est
alors pas surjectif. Il est toujours injectif, car $\widetilde{G}$ a une
section sur $\mathrm{Bor}$, la section unité. Enfin, $\widetilde{G}\to
\mathrm{Bor}$ étant affine et muni d'une section, un faisceau inversible sur
$\mathrm{Bor}$ est ample si et seulement si son image inverse l'est —
« indépendamment de l'hypothèse normal sur $S$ ! », note la marge.

\pagerange{12}{13}
\subsection*{Remarques (feuillets 12 et 13) : aucune résolution sans changement de base}

\textbf{a)} Si $T$ peut être tordu et que $S$ n'a qu'un nombre fini de
composantes connexes, le conoyau de l'injection
$\operatorname{Pic}(\mathrm{Bor})\to\operatorname{Pic}(\widetilde{G})$ est
fini. Pour $S$ intègre et normal, de point générique $\xi$, on a une suite
exacte
\[
  0\to\operatorname{Pic}(\mathrm{Bor})\to\operatorname{Pic}(\widetilde{G})
  \to H^{1}(\Gamma,X^{*}(T)_{\bar\xi}),
\]
fondée sur $\operatorname{Pic}(\widetilde{G}_{\eta})\simeq
\operatorname{Pic}(T_{\eta})\simeq\operatorname{Pic}(T_{\xi})$, $\eta$ le point
générique de $\mathrm{Bor}$. Si $S$ est régulier, on peut compléter la suite
par $\to 0$ ; la condition $H^{1}(\Gamma,X^{*}(T)_{\bar\xi})=0$ est alors
nécessaire et suffisante pour la surjectivité.\footnote{Le passage est criblé
de mots illisibles ; l'édition donne la structure que les formules
assurent.}

\textbf{b)} L'injectivité de
$\operatorname{Pic}(\mathrm{Bor})\to\operatorname{Pic}(\widetilde{G})$ et la
caractérisation des faisceaux amples restent vraies sans hypothèse de
normalité sur $S$.

\textbf{c)} Supposons $S$ normal et $H^{1}(\Gamma,X^{*}(T))=0$. Alors
$\operatorname{Pic}(S)\xrightarrow{\sim}\operatorname{Pic}(G')$.
L'injectivité vient de la section unité de $G'$ sur $S$. Soit
$\mathcal{L}\in\operatorname{Pic}(G')$ trivial sur la section unité. On a
\[
  \operatorname{Pic}(G')\hookrightarrow\operatorname{Pic}(G'^{\mathrm{reg}})
  \simeq\operatorname{Pic}(\widetilde{G}^{\mathrm{reg}})
  \simeq\operatorname{Pic}(\widetilde{G}),
\]
la première flèche injective par profondeur ; il suffit donc que l'image
inverse $p^{*}\mathcal{L}$ soit nulle. Comme
$\operatorname{Pic}(\widetilde{G})=\operatorname{Pic}(\mathrm{Bor})$, il suffit
de voir que sa restriction à la section unité $\mathrm{Bor}\to\widetilde{G}$
est nulle ; or $\mathrm{Bor}\to\widetilde{G}\to G'$ se factorise par la
section unité de $G'$ sur $S$, où $\mathcal{L}$ est trivial.\footnote{Le
feuillet note $\xi$ ce faisceau, lettre déjà prise par le point générique de
a) ; l'édition écrit $\mathcal{L}$. L'hypothèse $H^{1}=0$ est de l'édition :
l'argument se sert de
$\operatorname{Pic}(\widetilde{G})=\operatorname{Pic}(\mathrm{Bor})$, et le
feuillet 13 le reconnaît en donnant aussitôt la forme générale.} Sans
hypothèse sur $T$, le même diagramme

\begin{tikzcd}
G'^{\mathrm{reg}} \arrow[d, hook] & \widetilde{G}^{\mathrm{reg}} \arrow[l] \arrow[d, hook] \\
G' \arrow[d] & \widetilde{G} \arrow[l] \\
S \arrow[u, bend left=40] & \mathrm{Bor} \arrow[u, hook] \arrow[ul] \arrow[l]
\end{tikzcd}

donne une suite exacte $\operatorname{Pic}(S)\to\operatorname{Pic}(G')\to
H^{1}(\xi,X^{*}(T))$, de dernier terme fini.

\textbf{d)} Pour $S$ normal intègre, de point générique $\xi$, on a
$0\to\operatorname{Pic}(S)\to\operatorname{Pic}(G)\to
\operatorname{Pic}(G_{\xi})$, injective grâce à la section unité, et
$\operatorname{Pic}(G_{\xi})$ est fini, nul si $G_{\xi}$ est simplement
connexe.\footnote{Le feuillet dit « nul si $G_{\xi}$ est simplement connexe et
déployé » ; la nullité vaut déjà pour un groupe semi-simple simplement connexe
non déployé, mais la condition plus forte de la page n'est pas fausse.}
D'autre part
$\operatorname{Pic}(G)\to\operatorname{Pic}(G^{\mathrm{reg}})$ est injectif
sans hypothèse de normalité, par profondeur, le lieu non régulier étant de
codimension $\geqslant 3$ fibre par fibre. Si $S$ est régulier, cette flèche
est bijective, et le conoyau de
$\operatorname{Pic}(S)\to\operatorname{Pic}(G^{\mathrm{reg}})$ est
fini.\footnote{La page tire la finitude du conoyau de la seule injectivité ;
il y faut aussi la surjectivité, que donne la factorialité de $G$ quand $S$ est
régulier. D'où l'hypothèse ajoutée.}

Voici la conclusion, qui est le point du chapitre. Une résolution
$\widetilde{X}\to G$ des singularités de $f$ qui serait analogue à celle de
$f'$ — un isomorphisme au-dessus de $G^{\mathrm{reg}}$, avec un lieu
exceptionnel de codimension $\geqslant 2$ fibre par fibre — aurait un groupe
de Picard égal à celui de $G^{\mathrm{reg}}$, donc de rang nul sur
$\operatorname{Pic}(S)$ ; alors que celui de $\widetilde{G}$ contient le réseau
des poids, de rang $r$, avec son cône ample. Il n'y a donc pas de résolution
simultanée de $f$ sans toucher à $G^{\mathrm{reg}}$ : le changement de base
$T\to T/W$ est nécessaire.\footnote{L'énoncé du feuillet est plus bref :
« Cela montre qu'on ne peut trouver une résolution $\widetilde{X}$ des
singularités de $f$ analogue à celle de $f'$, sans toucher à $G^{\mathrm{reg}}$
(l'image inverse de $G-G^{\mathrm{reg}}$ dans $\widetilde{X}$ étant de codim.\
$\geqslant 2$ fibre par fibre…) ». Que les résolutions simultanées n'existent
qu'après le changement de base par $W$ est devenu, avec Brieskorn et Slodowy,
un énoncé classique. Un « e) » isolé, biffé, figure en bas de page.}

\pagerange{14}{19}
\section*{Résumé des théorèmes de conjugaison de Kostant (feuillets 14 à 19)}

La chemise du feuillet 14 porte « Résumé th.\ de conjugaison Kostant + lettre à
Kostant » ; la lettre est aux feuillets 21--22. Les feuillets 15 à 19 sont un
tapuscrit de sa main, corrigé et complété à l'encre, intitulé « Résumé de
quelques résultats de Kostant (sous-groupes simples de rang 1) », avec en marge
les numéros de l'article de 1959 qu'il résume (3.4, 3.5, 3.6, 4.2,
9.3).\footnote{Les lettres gothiques, absentes de la machine, sont ajoutées à
la main.}

\pagerange{15}{16}
\subsection*{Théorème 1 : sous-groupes de type $\mathrm{SL}(2)$}

$G$ est semi-simple sur un corps $k$ de caractéristique nulle, $e\neq 0$
nilpotent dans $\mathfrak{g}$. Pour $x\in\mathfrak{g}$ semi-simple et
$k\in\mathbf{Z}$, on note $\mathfrak{g}(k)$ le sous-espace propre de
$\operatorname{ad}(x)$ pour la valeur propre $k$\footnote{La frappe laisse un
blanc après « sous-espace propre de » ; la main ne l'a pas rempli.
$\operatorname{ad}(x)$ est ce que le sens demande.} ; si
$x=\operatorname{Lie}(i)(1)$ pour un cocaractère
$i\colon\mathbf{G}_{m}\to G$, ce sont les composantes isotypiques de
$\mathfrak{g}$ sous $\mathbf{G}_{m}$. On note $(x_{0},e_{0},f_{0})$ la base
habituelle de $\mathfrak{sl}_{2}$, $i_{0}$ le cocaractère diagonal de
$\mathrm{SL}(2)$, et $\operatorname{Fil}^{k}\mathfrak{g}=
\sum_{j\geqslant k}\mathfrak{g}(j)$.

\textbf{(i)} Les conditions suivantes sont équivalentes :

\begin{enumerate}
\item[a)] $[x,e]=2e$, et $\operatorname{ad}(e)^{k}\colon\mathfrak{g}(-k)\to
  \mathfrak{g}(k)$ est un isomorphisme pour tout $k\geqslant 1$ ;
\item[b)] $e\in\mathfrak{g}(2)$,
  $\operatorname{ad}(e)\colon\mathfrak{g}(0)\to\mathfrak{g}(2)$ est
  surjectif, et $x$ est un \emph{élément de Dynkin}, de la forme
  $\varphi(x_{0})$ pour un $\varphi\colon\mathrm{SL}(2)\to G$ ;
\item[c)] $e\in\mathfrak{g}(2)$ et $x\in[e,\mathfrak{g}]$ ;
\item[c')] $e\in\mathfrak{g}(2)$ et $x\in[e,\mathfrak{g}(-2)]$ ;
\item[d)] il existe un homomorphisme $\varphi\colon\mathrm{SL}(2)\to G$ tel que
  $d\varphi(e_{0})=e$ et $d\varphi(x_{0})=x$.
\end{enumerate}

Ces conditions entraînent que $x=\operatorname{Lie}(i)(1)$ pour un cocaractère
$i$ — en particulier que $x$ est semi-simple — et le $\varphi$ de d) est
unique.\footnote{Le tapuscrit écrit $f$ pour l'homomorphisme
$\mathrm{SL}(2)\to G$ ; l'édition écrit $\varphi$, pour ne pas le confondre avec
le morphisme de Steinberg ni avec le troisième vecteur d'un
$\mathfrak{sl}_{2}$-triplet. Dans b), il porte
$\operatorname{ad}(e)\colon\mathfrak{g}(0)\to\mathfrak{g}(-2)$ : $\operatorname{ad}(e)$
élève le degré de $2$, et (iii) au feuillet 16 écrit bien
$\mathfrak{g}(0)\to\mathfrak{g}(2)$ ; l'édition corrige. La seconde moitié de
b) est ajoutée à la main, et une ligne b') est barrée. L'équivalence de c) et
d) est le lemme de Morozov, celle de a) et d) la théorie des représentations
de $\mathfrak{sl}_{2}$.}

\textbf{(ii)} Pour tout nilpotent $e$, il existe $\varphi\colon\mathrm{SL}(2)\to
G$ avec $d\varphi(e_{0})=e$, ou, ce qui revient au même, avec
$\varphi(u_{0})=u$, où $u=\exp(e)$ et $u_{0}=\exp(e_{0})$ — c'est le théorème
de Jacobson--Morozov. Pour $e\neq 0$, ces $\varphi$ sont déterminés par
$x=d\varphi(x_{0})$, ils sont conjugués sous $Z=\operatorname{Cent}(e)=
\operatorname{Cent}(u)$, et ils forment un torseur sous $Z^{+}=Z\cap U$,
groupe unipotent, où $U$ est le groupe défini en (iv). L'ensemble des $x$
correspondants est l'espace affine
$x+\bigl(\mathfrak{g}^{e}\cap\operatorname{Fil}^{1}\mathfrak{g}\bigr)$, torseur
sous $\operatorname{Lie}(Z^{+})=\mathfrak{g}^{e}\cap
\operatorname{Fil}^{1}\mathfrak{g}$ pour l'addition. On a
$Z=Z_{0}\ltimes Z^{+}$, où $Z_{0}=Z\cap L=\operatorname{Cent}(\varphi)$ est
réductif.\footnote{Les deux définitions de $Z^{+}$ et de son algèbre de Lie
sont manuscrites, à la place de définitions dactylographiées biffées ; la
décomposition $Z=Z_{0}\cdot Z^{+}$ est une note marginale.}

\textbf{(iii)} Soit $x\neq 0$ un élément de Dynkin, de cocaractère $i$. Les $e$
tels que $(x,e)$ satisfasse aux conditions de (i) forment l'ouvert dense
$\mathfrak{g}(2)^{*}$ de $\mathfrak{g}(2)$ où
$\operatorname{ad}(e)\colon\mathfrak{g}(0)\to\mathfrak{g}(2)$ est surjectif ;
c'est un espace homogène — le tapuscrit corrige « torseur » — sous le groupe
réductif $L=\operatorname{Cent}(x)=\operatorname{Cent}(i)$.

\textbf{(iv)} Soit $x=\operatorname{Lie}(i)(1)$ un élément de Dynkin, et $P$,
$U$, $L$ les sous-groupes d'algèbres de Lie
$\sum_{k\geqslant 0}\mathfrak{g}(k)$, $\sum_{k\geqslant 1}\mathfrak{g}(k)$ et
$\mathfrak{g}(0)$. Alors $P$ est parabolique, $U$ son radical unipotent, $L$ un
sous-groupe de Levi. Si $u=\exp(e)$ est adapté à $i$, $P$ et $U$ ne dépendent
que de $u$ (deux $x$ associés au même $u$ donnent le même $P$), de même que
$H=i(\mathbf{G}_{m})\cdot U$ et que $i$ modulo $U$. L'application
$g\mapsto\operatorname{Ad}(g)x$ est un isomorphisme de $U$ sur
$x+\operatorname{Fil}^{1}\mathfrak{g}$, et $g\mapsto\operatorname{Ad}(g)e$
applique $U$ sur $e+\operatorname{Fil}^{3}\mathfrak{g}$, et $P$ sur
$\mathfrak{g}(2)^{*}+\operatorname{Fil}^{3}\mathfrak{g}$.\footnote{C'est ce
qu'on appelle aujourd'hui le parabolique de Jacobson--Morozov de $e$ ; les
valeurs de $\alpha(i)$ sur les racines simples forment le diagramme de Dynkin
pondéré de l'orbite de $e$.}

\pagerange{17}{19}
\subsection*{Interprétation par les filtrations du foncteur fibre}

Le tapuscrit traduit ensuite la théorie en termes de filtrations du foncteur
oubli sur la catégorie des représentations de dimension finie de $G$ —
filtrations compatibles au produit tensoriel et exactes, ce qui est
automatique en caractéristique nulle.\footnote{C'est le cadre que développera
la thèse de Saavedra Rivano, \emph{Catégories tannakiennes} (Lecture Notes 265,
1972), sous sa direction ; le tapuscrit en est antérieur ou contemporain.}

Une classe de conjugaison $c$ de telles filtrations, rationnelle sur $k$
(sans qu'elle ait nécessairement un représentant rationnel), revient à un
élément de $X_{*}(T_{0})/W$ invariant par l'image $\Gamma$ de
$\operatorname{Gal}(\bar k/k)$ dans le groupe des automorphismes extérieurs du
groupe de Chevalley $G_{0}$ associé — autrement dit à un cocaractère dominant
stable par $\Gamma$. On a alors :

\begin{itemize}
\item le schéma $\underline{\varnothing}_{c}$ des filtrations de classe $c$,
  espace homogène sous $G$ dont les stabilisateurs sont des paraboliques ; il
  s'identifie à l'ouvert et fermé $\mathrm{Par}_{c'}$ de
  $\underline{\mathrm{Par}}(G)$ formé des paraboliques d'un type $c'$ ;
\item le schéma $\mathrm{ST}_{c}$ des cocaractères
  $i\colon\mathbf{G}_{m}\to G$ de classe $c$, orbite rationnelle dans
  $\underline{\mathrm{Hom}}_{\mathrm{gr}}(\mathbf{G}_{m},G)$, isomorphe au
  schéma $(\mathrm{CL})_{c'}$ des couples de Levi $(P,L)$ avec $P$ de type
  $c'$ : à $i$ on associe $P$ et $\operatorname{Cent}(i)$ ;
\item le morphisme $\mathrm{ST}_{c}\to\underline{\varnothing}_{c}$ (en termes de
  filtrations), qui s'identifie à $(\mathrm{CL})_{c'}\to\mathrm{Par}_{c'}$ (en
  termes de paraboliques).
\end{itemize}

En caractéristique nulle (« j'ignore dans quelle mesure c'est vraiment
nécessaire »), soit $\mathrm{Dyn}_{c}$ le schéma des
$\varphi\colon\mathrm{SL}(2)\to G$ tels que $\varphi\circ i_{0}$ soit de classe
$c$. S'il est non vide et $c\neq 1$ — $c$ est alors une \emph{classe de
Dynkin} —, c'est un espace homogène sous $G$, de stabilisateur
$Z_{0}$.\footnote{La main ajoute : « il faudrait regarder la structure de ces
$Z_{0}$ ; dimension, semi-simplicité ». Le tapuscrit biffe une parenthèse qui
identifiait $\mathrm{Dyn}_{c}$ au torseur $\underline{\mathrm{Isom}}(G_{1},G)$,
$G_{1}$ la forme quasi-déployée.} Soit $\mathrm{Un}_{c}$ le schéma des
unipotents de la classe qui correspond à $c$. Les morphismes $\varphi\mapsto
\varphi\circ i_{0}$ et $\varphi\mapsto\varphi(u_{0})$, et un morphisme
$\mathrm{Un}_{c}\to\underline{\varnothing}_{c}$ que « Deligne nous fournit »,
forment un carré commutatif d'espaces homogènes

\begin{tikzcd}
\mathrm{ST}_{c}\simeq(\mathrm{CL})_{c'} \arrow[d] & \mathrm{Dyn}_{c} \arrow[l] \arrow[d] \\
\underline{\varnothing}_{c}\simeq\mathrm{Par}_{c'} & \mathrm{Un}_{c} \arrow[l]
\end{tikzcd}

de stabilisateurs $L$, $Z_{0}$, $P$ et $\operatorname{Cent}(e)$.\footnote{Le
carré des stabilisateurs est dessiné à la main, à côté du premier. Le morphisme
de Deligne associe à un unipotent la filtration canonique qu'il définit ; le
rapprocher de la « filtration de monodromie » de \emph{La conjecture de Weil
II} (1980), 1.6, est une remarque de l'édition, non du dossier.}

\textbf{Points rationnels.} $\mathrm{Par}_{c'}$ a un point rationnel si et
seulement si $\mathrm{Dyn}_{c}$ en a un. Dans un sens, un $\varphi$ rationnel
définit $i=\varphi\circ i_{0}$, donc $P$. Dans l'autre, $P$ rationnel étant
donné, on en choisit un Levi, c'est-à-dire un $i$ de classe $c$ qui le définit,
et les $e$ adaptés à $i$ forment l'ouvert dense $\mathfrak{g}(2)^{*}$ de
$\mathfrak{g}(2)$, qui a un point rationnel. Dans le cas principal, où
$Z_{0}$ est trivial et $P$ un Borel, cela dit qu'un sous-groupe principal
$\mathrm{SL}(2)$ rationnel existe si et seulement si $G$ est
quasi-déployé.\footnote{Le tapuscrit en tire que « pour certaines classes de
paraboliques (qu'on pourrait appeler les paraboliques de Dynkin-Kostant)
l'existence d'un parabolique de cette classe rationnel sur $k$ implique déjà
que $G$ est quasi-déployé ». L'argument donné ne l'établit que dans le cas
principal, où $P$ est un Borel et où la conclusion est immédiate ; pour une
autre classe de Dynkin, un $\varphi$ rationnel n'entraîne que l'isotropie.
L'édition s'en tient à ce qui est démontré.} L'argument vaut sur un anneau
semi-local de caractéristique nulle.

Quand $G$ n'est pas une forme intérieure, il faut vérifier que la classe $c'$
de $P$ provient d'une classe de Dynkin $c$ \emph{invariante par} $\Gamma$, ce
qui ne résulte pas de l'invariance de $c'$ tant qu'on ne sait pas si $c\mapsto
c'$ est injectif sur les classes de Dynkin. Il ne l'est pas : une note marginale
le dit pour $A_{2}$ et $A_{4}$.\footnote{Lecture partiellement incertaine :
« Il n'est en effet pas injectif pour $A_{2}$, $A_{4}$ (cela vaut pour
$A_{3}$) ». Le calcul sur les diagrammes pondérés confirme : dans $A_{2}$,
l'orbite régulière ($22$) et l'orbite minimale ($11$) ont toutes deux un Borel
pour parabolique ; dans $A_{4}$, la régulière et la sous-régulière ($2112$) de
même ; dans $A_{3}$, les orbites $[3,1]$ ($202$) et $[2,1,1]$ ($101$) ont le
même parabolique, ce qui justifie la parenthèse.} D'où l'intérêt, pour
l'arithmétique, de déterminer les classes de paraboliques de Dynkin--Kostant des
groupes simples, et de comparer leur stabilisateur dans le groupe des
automorphismes extérieurs à celui des $c$ correspondants.

Si une filtration du foncteur oubli est définie par un unipotent, la classe de
celui-ci (ou $c$) est déterminée. Les unipotents qui la définissent sont les
$\exp(e)$, $e\in\mathfrak{g}(2)^{*}+\operatorname{Fil}^{3}\mathfrak{g}$ ; ils
forment une seule classe de conjugaison sous $P$ — non sous $U$, dont les
classes sont les $\exp(a+\operatorname{Fil}^{3}\mathfrak{g})$,
$a\in\mathfrak{g}(2)^{*}$.

\textbf{Cas principal.} La classe de $\varphi$ est \emph{principale} au sens de
Dynkin--Kostant lorsque $e$ est régulier, ce qui équivaut à ce que
$\mathfrak{g}(0)$ soit une sous-algèbre de Cartan ; on a alors
$\mathfrak{g}(1)=0$. Le parabolique associé est un Borel — la réciproque est
fausse, comme on vient de le voir —, et $Z_{0}$ est trivial pour $G$ adjoint :
$\mathrm{Dyn}_{c}$ est un torseur sous $G$.\footnote{Le tapuscrit écrit :
« Pour qu'on ait [blanc]$(1)=0$, il faut et suffit que la classe … soit
“principale” ». Le blanc n'est pas rempli. Lu, comme le propose la
transcription, $\mathfrak{g}(1)=0$, l'équivalence est fausse : $\mathfrak{g}(1)=0$
caractérise les classes \emph{paires}, dont la principale n'est qu'une (dans
$\mathfrak{sl}_{4}$, l'orbite $[2,2]$ est paire et non principale).
L'édition donne la caractérisation correcte et garde l'implication. La marge
note : « ex. Cas $A_{2}$, $A_{4}$ : les paraboliques associés aux unipotents
“sous-réguliers” sont des Borels ! »}

\pagerange{21}{22}
\section*{La lettre à Kostant (Massy, 22 octobre 1969)}

Lettre dactylographiée, en anglais, signée de sa main.\footnote{Le titre est de
l'édition. Quelques mots sont ajoutés à la main — les deux carrés de
$\operatorname{card}(W)^{2}$, « of $\widetilde{G}$ », « generating », « $n-$ »,
« all » — et une note de bas de page. La lettre est de lui ; la transcription
la donne en entier, et l'édition la suit ici en français.} Elle dit l'occasion
de la relecture : un « joli travail » de Brieskorn sur les singularités de
Klein a ramené Grothendieck aux deux articles de Kostant, qui « appellent une
reconsidération soigneuse dans le cadre de la géométrie algébrique ». Suivent
des questions.

\textbf{Quadruplets rigidifiants.} Pour $G$ simple adjoint complexe, classer les
quadruplets $(T,B,T',B')$ où $T$ et $T'$ sont des tores maximaux « en
apposition », au sens de Kostant, et $B\supset T$, $B'\supset T'$ des Borels.
Le nombre de classes de conjugaison est annoncé égal à
\[
  \frac{\operatorname{card}(W)^{2}}{h\,z\,\varphi(h)} ,
\]
et à $\operatorname{card}(W)^{2}/hz$ si l'on se donne en outre un générateur du
groupe d'ordre $h$, $T\cap N(T')$. Le stabilisateur d'un tel quadruplet est
$T\cap T'=e$ : la donnée rend $G$ entièrement « rigide », aux automorphismes
extérieurs près.\footnote{Note de sa main au pied du feuillet :
$N(T)\cap N(T')$ est une extension de $(\mathbf{Z}/h\mathbf{Z})^{*}$ par
$\mathbf{Z}/h\mathbf{Z}\times\mathfrak{z}$, donc d'ordre $hz\varphi(h)$ ; le
dénombrement en découle si ce groupe agit librement sur les
$\operatorname{card}(W)^{2}$ choix de $(B,B')$. Le dossier ne démontre ni
l'une ni l'autre affirmation. Les nombres sont entiers dans les cas qu'on
vérifie à la main : $1$ pour $A_{1}$, $2$ pour $A_{2}$, $4$ pour $B_{2}$, $12$
pour $G_{2}$.} Questions : $B$ et $B'$ sont-ils en position générale,
c'est-à-dire $B\cap B'$ est-il un tore maximal, pour certains ou pour tous ces
quadruplets ? Y a-t-il une classe \emph{distinguée}, invariante par les
automorphismes extérieurs ? Quels sont les groupes d'automorphismes de ces
quadruplets — isomorphes, dit-il, à des sous-groupes du groupe des
automorphismes extérieurs ?

\textbf{Le rang de l'application de Kostant.} Dans une autre direction : comment
varie le rang de la différentielle de $\mathfrak{g}\to\mathfrak{h}/W$ (et de
celle de Steinberg $G\to T/W$) ? Il y a des raisons de croire que sur les
orbites unipotentes de dimension $n-r-2$, juste en dessous du maximum $n-r$, le
rang est $r-1$ quand toutes les racines ont même longueur — types $A$, $D$,
$E_{6}$, $E_{7}$, $E_{8}$, ceux des singularités de Klein — et $<r-1$ dans les
autres cas ; et que dans le premier cas il n'y a qu'une orbite de cette
dimension, propriété peut-être caractéristique du cas « homogène ».\footnote{Il
y a en fait, dans tout groupe simple, une seule orbite nilpotente de dimension
$n-r-2$, l'orbite sous-régulière ; l'unicité n'est donc pas propre au cas des
racines de même longueur. Le lien avec les singularités de Klein que la lettre
devine est devenu le théorème de Brieskorn (congrès de Nice, 1970), conjecturé
par Grothendieck : une tranche transverse à l'orbite sous-régulière coupe le
cône nilpotent selon une singularité simple de même type, et la restriction de
l'application de Kostant à la tranche en est la déformation semi-universelle
dans le cas $ADE$ (Slodowy, 1980, pour les autres types). L'édition ne tranche
pas l'énoncé sur le rang dans les types non simplement lacés.}

\textbf{Les fibres de Springer.} Quelle est la dimension de la variété des
Borels qui contiennent un élément $g\in G$ (resp.\ dont l'algèbre de Lie
contient $x\in\mathfrak{g}$) ? Si l'orbite est de dimension $2N-2i$, il
s'attend à la dimension $i$ ; et l'inégalité $\leqslant i$ entraînerait, en
généralisant un argument de profondeur de Brieskorn, que les singularités des
fibres de $\mathfrak{g}\to\mathfrak{h}/W$ (resp.\ de $G\to T/W$ pour $G$
simplement connexe) sont \emph{rationnelles} : pour une résolution $F'\to F$
d'une fibre, $R^{i}f_{*}\mathcal{O}_{F'}=0$ pour $i>0$.\footnote{La première
attente est devenue la formule $\dim\mathcal{B}_{x}=\frac12(\dim
Z_{G}(x)-r)$, établie au milieu des années 1970 (Steinberg, Spaltenstein) ;
la seconde, la rationalité des singularités, l'a été à la même époque pour le
cône nilpotent (Hesselink, 1976). C'est la question g) du feuillet 6.}

La lettre se clôt sur une phrase qui dit l'attitude : « Peut-être la réponse à
la plupart des questions que je pose est-elle bien connue, et suis-je seulement
ignorant ; pour certaines d'entre elles, il serait en vérité scandaleux que les
experts ne connaissent pas la réponse ! »

\pagerange{23}{35}
\section*{Diagrammes d'espaces homogènes : éléments de Kostant, de Coxeter, points cycliques (feuillets 23 à 35)}

Le titre est celui de la chemise du feuillet 23. Les notes qui suivent sont sur
papier libre ou au verso de feuillets dactylographiés étrangers au dossier
(27, 29, 31, 35), qu'on ne fait que signaler. Elles relèvent dans $G$ les
objets que les feuillets 1--2 décrivaient dans $W$, sur un corps $k$ de
caractéristique nulle ou première à $h$.

\pagerange{25}{26}
\subsection*{Le tableau des notions (feuillets 25 et 26)}

Le feuillet 25 dresse, en trois colonnes, une liste de notions, leurs sigles,
et le groupe qui les stabilise.

\begin{itemize}
\item \emph{Éléments de Kostant} $\kappa\colon\mu_{h}\to G$ et \emph{éléments de
  Coxeter} $\gamma\in G(k)$ (E Kost, E Cox) : équivalents en caractéristique
  nulle dès qu'on s'est donné $\omega\in\mu_{h}^{*}(k)$, par $\gamma=\kappa(\omega)$ ;
  stabilisateur $\mathfrak{z}\cdot T$, produit semi-direct.
\item \emph{Sous-groupes} de Kostant et de Coxeter $\mu\hookrightarrow G$
  (SG Kost, SG Cox) ; stabilisateur extension de
  $(\mathbf{Z}/h\mathbf{Z})^{*}$ par $\mathfrak{z}\cdot T$.
\item \emph{Couples} de Kostant $(\kappa,T')$, ou $(\kappa,D)$ avec $D$ une
  droite régulière de $\mathfrak{g}(1)$, et de Coxeter $(\gamma,T)$
  (C Kost, C Cox) ; stabilisateur $\mu_{h}\times\mathfrak{z}$.
\item \emph{Précouples} $(\mu\subset G,T')$, c'est-à-dire couples de tores en
  apposition (Toapp) ; stabilisateur extension de
  $(\mathbf{Z}/h\mathbf{Z})^{*}$ par $\mu_{h}\times\mathfrak{z}$.
\item \emph{Droites cycliques} $D$ (D Cy) ; stabilisateur $\mu_{h}\cdot T'$.
\item \emph{Couples de Kostant--Killing} $((B,T),T')$, équivalemment
  $(B,T,D)$ ; stabilisateur $\mu_{h}$.
\item \emph{Couples de Kostant normalisés} (C Kost norma) ; stabilisateur
  $\mathfrak{z}$.
\item \emph{Couples de Kostant--Killing normalisés} $((T,B),x)$
  (Kost Kill norma) ; stabilisateur trivial.
\end{itemize}

Une marge précise que ces équivalences ne valent qu'en caractéristique nulle
« (ou car.\ $p>h$ ?) », et une note définit la normalisation : le choix d'un
élément d'une droite canonique $\Delta$, commune aux droites cycliques modulo
$\mu_{h}$, choix privilégié au signe près.\footnote{Les sigles ont été récrits
plusieurs fois et les biffures y sont lourdes ; plusieurs mots de la note sont
incertains. L'édition donne les lignes dans la dernière forme lisible.}

Le feuillet 26 range les mêmes données en deux treillis — celui des
stabilisateurs, de $0$ à $N(T)$ et $N(T')$ en passant par $\mathfrak{z}$,
$\mu$, $\mu\times\mathfrak{z}$, $\mathfrak{z}\cdot T$, $\mu\cdot T'$, $N(\mu)$,
$N(T)\cap N(T')$ ; et celui des espaces, de « Kost-Kill-norma » en haut jusqu'aux
tores maximaux en bas — avec deux renvois : certains isomorphismes ne sont
définis que sur $\mathring{\mu}_{h}^{*}$, d'autres que si $(h,p)=1$. Une
dernière ligne note que l'existence de ces objets dans le cas quasi-déployé
s'obtient par descente.\footnote{L'édition ne redessine pas ces deux
diagrammes, que la transcription donne trait pour trait ; le long trait qui
descend de C Kos y est de lecture incertaine.}

\pagerange{28}{30}
\subsection*{Les groupes d'une paire de tores en apposition (feuillets 28 et 30)}

Le feuillet 28 identifie les couples de tores en apposition $(T,T')$ aux
sous-groupes de Coxeter $\mu$, et le choix d'un isomorphisme $\mu\simeq\mu_{h}$
au choix d'une droite cyclique $D$. Puis il croise, dans une grille, les
sous-groupes $T$, $C(\mu)$ (centralisateur de $\mu$, que la page note, comme le centre, d'un 3 bouclé), $N(\mu)$, $N(T)$
avec leurs intersections par $N(T')$, par $N(D)$ et par $T'$, en notant sur chaque
trait le quotient ($\mathfrak{z}$, $(\mathbf{Z}/h\mathbf{Z})^{*}$,
$T/\mu$, $\mu$, $T'/\mathfrak{z}$) et en marquant les carrés cocartésiens. Il
en retient trois égalités :
\[
  N(T)\cap N(D)=C(\mu)\cap N(D),\qquad
  N(T)\cap T'=C(\mu)\cap T',\qquad
  N(T)\cap N(T')=N(\mu)\cap N(T'),
\]
et une extension canonique de $(\mathbf{Z}/h\mathbf{Z})^{*}$ par
$\mu\times\mathfrak{z}$, dont les deux extensions induites, par $\mu$ et par
$\mathfrak{z}$, se réalisent dans $W'$ et dans $W$.\footnote{C'est la question
de la marge du feuillet 1, reprise au niveau de $G$. La fin du paragraphe ne
se lit que par mots isolés.}

Le feuillet 30 refait la grille pour un élément de Coxeter $\gamma$ de
$N(T)$. Il y lit $T\cap T'=0$ et
\[
  T'^{\gamma}\simeq P'/(1-c)P'(1)\simeq\mathfrak{z},
\]
où $P'$ est le réseau des cocaractères de $T'$ (copoids, $G$ étant adjoint) et
$c$ l'élément de Coxeter qu'induit $\gamma$ : comme $1-c$ est inversible sur
$\mathbf{Q}$, les points fixes de $\gamma$ dans $T'$ forment le groupe fini
$P'/(1-c)P'$, tordu à la Tate, et $(1-c)$ envoie le réseau des copoids sur celui
des coracines, d'où un groupe isomorphe à $P^{\vee}/Q^{\vee}$, d'ordre
$z$.\footnote{L'explication de l'isomorphisme est de l'édition ; la page ne
donne que la formule, sous la grille.} Deux justifications suivent : tout
élément de $T$ qui normalise $D$ est neutre ; tout élément de $W$ qui
centralise $c$ en est une puissance (le corollaire 1 des feuillets 1--2).

\emph{À prouver} : tout élément de $T$ qui normalise $D$ est dans
$K(\mu_{h})$. La page amorce un calcul d'indice, le déterminant
\[
  \det\begin{pmatrix}
    n_{1}+1 & n_{2} & \cdots & n_{r}\\
    n_{1} & n_{2}+1 & \cdots & n_{r}\\
    \vdots & & \ddots & \vdots\\
    n_{1} & n_{2} & \cdots & n_{r}+1
  \end{pmatrix}=1+\sum_{i}n_{i},
\]
qui vaut $h$ quand les $n_{i}$ sont les coefficients de la plus grande racine.
Elle conclut que $N(T)\cap N(D)=N(\mu)\cap N(D)$, puisque $\mu\simeq\mu_{h}$
est formé des éléments de $T$ qui normalisent $D$.\footnote{L'énoncé à
prouver est établi au feuillet 46, (iii) c), sous la forme : le stabilisateur
dans $T$ d'une droite cyclique est $K(\mu_{h})$. La matrice est esquissée sur la
page ; sa valeur est juste (lemme du déterminant d'une perturbation de rang
un).}

\pagerange{33}{33}
\subsection*{Rappels de Kostant (feuillet 33)}

\textbf{a)} Pour $\mu_{h}\hookrightarrow T\hookrightarrow G$ fixé, les droites
cycliques $D$ normalisées par $\mu_{h}$ forment un espace homogène sous $T$, de
stabilisateur $\mu_{h}$. Elles correspondent bijectivement aux tores maximaux
$T'$ en apposition avec $T$ par $\mu_{h}$ — normalisés par $\mu_{h}$, qui
agit sur $\operatorname{Lie}T'$ avec le caractère $1$ pour valeur propre — : à
$D$ on associe son centralisateur, à $T'$ la droite propre correspondante de
$\operatorname{Lie}T'$. On a $T\cap N(T')=\mu_{h}$.

\textbf{b)} Les droites cycliques de $\mathfrak{g}$, modulo $\mu_{h}^{*}$,
s'identifient canoniquement à une droite $\Delta$ définie sur $k$, obtenue par
descente à partir de $D^{\otimes h}$, ou comme l'image des lieux cycliques par
le morphisme de Kostant $\mathfrak{g}\to\mathfrak{t}/W$. On a
$\mathfrak{g}_{\psi}\simeq\Delta$, ou $\Delta^{-1}$, car
$\Delta^{\otimes 2}\simeq\mathbf{1}$.\footnote{La plus grande partie de b) et
de la note oblique qui l'accompagne est rapide et effacée ; l'édition ne
reconstitue pas les liaisons. La note parle d'un torseur et d'une obstruction
de Kostant dans un groupe $H^{2}(Q;E;\mu_{h})$, nulle si un certain
$H^{2}(Q,\mu_{h})$ l'est ; ce qu'elle obstrue ne se lit pas.}

\pagerange{34}{34}
\subsection*{Les rigidifications de Kostant (feuillet 34)}

Soit, sur une base où $(h,p)=1$, le schéma $\mathrm{RigKos}$ des
\emph{rigidifications de Kostant} $(T,T',B,B')$ : $T$ et $T'$ tores maximaux en
apposition, $B\supset T$ et $B'\supset T'$ des Borels. C'est un schéma lisse
sur $S$, sur lequel $G$ opère librement, et $\mathrm{RigKos}/G$ est fini étale
de rang $\operatorname{card}(W)^{2}/\varphi(h)hz$ — le nombre de la lettre.
Le feuillet le décrit comme une représentation du groupe des automorphismes
extérieurs dans un schéma fini étale sur $\operatorname{Spec}\mathbf{Z}[1/h]$.

Dans le cas déployé, un épinglage fournit $\kappa\colon\mu_{h}\to T$ et la
droite $D$ engendrée par
\[
  e=e_{\alpha_{1}}+\cdots+e_{\alpha_{r}}+e_{-\psi},
\]
définie au signe près ; $T'$ est le centralisateur de $e$. Le schéma $Q$ des
Borels contenant $T$ est un $W$-torseur trivial ; celui, $Q'$, des Borels
contenant $T'$ est un $W'$-torseur « très tordu ». Le sous-groupe
$\Phi$, image de $N(T)\cap N(T')$ dans $W\times W'$, extension de
$(\mathbf{Z}/h\mathbf{Z})^{*}$ par $\mathfrak{z}\times\mu_{h}$, opère sur
$Q\times Q'$, et le quotient cherché est $(Q\times Q')/\Phi$ ; les
automorphismes extérieurs y opèrent par leur action $\pm 1$ sur
$\mathfrak{g}_{-\psi}$.\footnote{Le feuillet ne nomme pas $Q$, et la lettre
$\Phi$ est surchargée. « librement » est de lecture incertaine. Les quatre
dernières lignes, à l'encre brune, ne se lisent que par fragments. Rien de ce
feuillet n'est démontré.}

\pagerange{36}{39}
\section*{Liste d'espaces homogènes remarquables et leurs relations (feuillets 36 à 39)}

La chemise du feuillet 36 précise au crayon : « liés aux sous-groupes simples
principaux de rang 1 ». Le feuillet 37 dresse la liste des schémas en jeu :

\begin{itemize}
\item les sous-groupes simples de rang $1$ principaux (SSP), seuls ou munis
  d'un épinglage (SSPEP, contenu dans
  $\underline{\mathrm{Hom}}^{*}_{\mathrm{gr}}(\mathrm{SL}(2),G)$), d'un couple
  de Killing, d'un tore maximal, d'un Borel, d'un Borel épinglé ;
\item les tores de dimension $1$ principaux (STP) et les homomorphismes
  principaux $\mathbf{G}_{m}\to G$ ($\mathrm{ST}_{0}\mathrm{P}$) ;
\item les sous-groupes additifs principaux (SAP), pointés
  ($\mathrm{SAP}_{0}\subset\underline{\mathrm{Hom}}_{\mathrm{gr}}(\mathbf{G}_{a},G)$),
  et les unipotents principaux (EUP) — canoniquement isomorphes en
  caractéristique nulle ;
\item les éléments de Coxeter et de Kostant, isomorphes une fois choisie une
  racine primitive $h$-ième de l'unité ;
\item les droites cycliques (DC), les éléments cycliques, les éléments
  cycliques principaux ;
\item les couples de tores maximaux en apposition (CTAP), de Coxeter
  $(T,\gamma\in W(T))$ (CCox), de Kostant $(\kappa,D)$ (CKost), ces deux
  derniers isomorphes sur $\mu_{h}^{*}$.
\end{itemize}

Le feuillet 38 dessine le diagramme de ces espaces homogènes sous $G$, de
$\underline{\mathrm{Isom}}(G_{0},G)\simeq\mathrm{SSPEP}$ en haut à
$\mathrm{STP}$ et $\mathrm{SSP}$ en bas, avec
$\mathrm{Kill}\simeq\mathrm{ST}_{0}\mathrm{P}$ et
$\mathrm{SAP}_{0}\simeq\mathrm{EUP}$ ; puis celui des stabilisateurs
correspondants, avec les dimensions sur les traits. Les dimensions se
vérifient : en partant du sous-groupe principal $H$ (dimension $3$), de son
tore $t$ et de son sous-groupe unipotent $b_{u}$ (dimension $1$), on monte au
tore maximal $T$ et au centralisateur $Z(b_{u})$ (dimension $r$), puis au
normalisateur $N(b_{u})=Z(b_{u})\cdot t$ (dimension $r+1$).\footnote{Le
feuillet 38 est au verso d'une page dactylographiée étrangère ; le
feuillet 39, d'un tapuscrit sur les topos, ne porte que des ajouts à ce texte.
La transcription donne les deux diagrammes du feuillet 38 ; l'édition n'en
retient que la vérification des dimensions. Une note du feuillet 37 dit que les
espaces des éléments principaux « deviennent isomorphes une fois choisie une
racine primitive $h$-ième de $1$ (donc dans $\mu_{h}^{*}$ !) ».}

\pagerange{40}{41}
\section*{Résultats de Kostant sur les classes remarquables (feuillets 40 et 41)}

$G$ est semi-simple adjoint sur $\mathbf{C}$.

\textbf{a)} Les classes de conjugaison de sous-groupes simples de rang $1$
(ou de sous-groupes additifs) correspondent aux classes d'unipotents de $G$
(ou de nilpotents de $\mathfrak{g}$), et à certaines classes de cocaractères
$\mathbf{G}_{m}\to G$. Dans $\operatorname{Hom}(\mathbf{G}_{m},T)\simeq
\mathbf{Z}^{\Sigma}$, ces cocaractères « admissibles », pris dans la chambre
dominante $\mathbf{N}^{\Sigma}$, ont des composantes égales à $0$, $1$ ou $2$
— pas toutes les combinaisons : ce sont les diagrammes de Dynkin pondérés.

\textbf{b)} Le cas où toutes les composantes valent $2$ est admissible : il
correspond aux unipotents principaux (réguliers), et aux homomorphismes
principaux.\footnote{La page écrit « toutes les composantes sont $1$ ». Pour le
cocaractère de Dynkin d'un $\mathrm{SL}(2)$ principal, elles valent $2$ ; le
cocaractère à composantes $1$ est $\rho^{\vee}$, sa moitié, l'« homomorphisme
principal » du feuillet 43 (voir les conventions).} Puis une série de
centralisateurs et de normalisateurs : celui d'un unipotent principal est un
groupe unipotent commutatif de dimension $r$ ; celui d'un homomorphisme
principal $\mathbf{G}_{m}\to G$ est un tore maximal, d'où une bijection entre
ces homomorphismes et les couples de Killing ; le normalisateur d'un tore de
dimension $1$ principal est engendré par $T$ et un élément d'ordre $2$ ; le
stabilisateur d'un sous-groupe principal de rang $1$, muni d'un Borel ou d'un
couple de Killing, se lit dans ce sous-groupe. En marge : sur une base
quelconque, il existe un sous-groupe principal de rang $1$ déployé si et
seulement si $G$ est quasi-déployé.\footnote{Plusieurs de ces phrases sont
lacunaires ; le normalisateur d'un sous-groupe additif principal reste sans
description lisible. Une note marginale nomme l'élément $w_{0}$ qui envoie la
chambre sur son opposée, sans dire à quel énoncé elle se rattache ; c'est,
d'après le feuillet 43, l'élément qui normalise le tore principal.}

\textbf{c)} Soit $u\colon\mathbf{G}_{m}\to G$ un homomorphisme principal, $h$
le nombre de Coxeter (la somme des coefficients de la plus grande racine, plus
$1$) et $\zeta$ une racine primitive $h$-ième de l'unité. Les éléments
$u(\zeta)$ sont réguliers et tous conjugués ; leur image dans le groupe de Weyl
de leur tore centralisateur est un élément de Coxeter, et réciproquement les
représentants d'un élément de Coxeter de $W=N(T)/T$ sont réguliers. La marge
conseille de considérer plutôt la classe de l'homomorphisme
$\mu_{h}\to G$.\footnote{Le paragraphe c) est surchargé d'ajouts dont l'ordre
ne se reconstitue pas avec certitude ; l'édition donne ce que les segments
lisibles assurent. La phrase sur le centralisateur d'un élément régulier
principal, suivie d'un « (?) » entouré, est laissée en suspens par la page.}

\textbf{d)} Le groupe $W$ opérant sur $\mathfrak{t}$, il existe des droites
propres d'un élément de Coxeter de valeur propre une racine primitive $h$-ième,
qui sont régulières, et conjuguées sous $W$ ($G$ simple) : ce sont les droites
cycliques des feuillets 1--2. Au feuillet 41, la phrase passe de $\mathfrak{t}$ à
$\mathfrak{g}$ : les droites de $\mathfrak{g}$ propres sous un élément régulier
principal, de valeur propre une racine primitive $h$-ième, sont conjuguées ; on
les appelle \emph{cycliques}, et leurs points non nuls \emph{éléments
cycliques}. Pour un couple de Killing donné, on en construit en prenant
\[
  \sum_{i=1}^{r}\lambda_{i}e_{\alpha_{i}}+\lambda_{0}e_{-\psi},
  \qquad\text{tous les }\lambda\text{ non nuls}.
\]
Le centralisateur d'un élément cyclique est un tore maximal $T'$, et le
normalisateur de la droite $D$ qu'il engendre est un groupe $\widetilde{T}'$,
$T'\subset\widetilde{T}'\subset N(T')$, dont l'image dans le groupe de Weyl est
cyclique d'ordre $h$, engendrée par un élément de Coxeter.\footnote{La page
appelle $T$ ce tore centralisateur et $\mu$ le dernier coefficient ;
l'édition écrit $T'$ et $\lambda_{0}$, pour garder $T$ au tore du couple de
Killing et $\mu$ au sous-groupe de Coxeter. Les justifications sont
interlinéaires et leur place incertaine.} Pour une droite cyclique, $D/\mu_{h}$
s'identifie canoniquement à la droite standard — ce qui se vérifie dans un cas
quasi-déployé, où l'on a l'élément cyclique privilégié
$\sum_{i}e_{\alpha_{i}}+e_{-\psi}$ — ; les éléments qui y donnent $1$ sont les
\emph{éléments cycliques principaux}, définis au signe près. Ils forment un
espace homogène sous $G$, de stabilisateur un tore maximal.\footnote{La page
dit « torseur » avec un stabilisateur non trivial ; l'édition écrit « espace
homogène ». Une marge demande si l'on peut choisir un générateur canonique de
$e_{-\psi}$, c'est-à-dire si un épinglage identifie $\mathfrak{g}_{-\psi}$ à
$\mathfrak{g}_{\psi}$ — c'est l'ambiguïté de signe du feuillet 33 et du
feuillet 34.}

\textbf{e)} Considérons les couples $(T,\gamma)$ d'un tore maximal et d'un
$\gamma\in N(T)$ dont l'image dans $W$ est un élément de Coxeter. Deux tels
couples sont conjugués — et la page s'arrête là, la justification entre
crochets étant illisible.

\pagerange{42}{47}
\section*{Résumé de quelques résultats de Kostant (feuillets 42 à 47)}

Seconde mouture, manuscrite et mise au propre, d'un résumé de Kostant ; elle ne
reprend pas le tapuscrit des feuillets 15--19, mais les objets des
feuillets 21--41, sur une base. $G$ est semi-simple adjoint sur $S$, de type
constant.

\pagerange{43}{44}
\subsection*{(i) Homomorphismes principaux}

À tout couple de Killing $(T,B)$, on associe un cocaractère
$i\colon\mathbf{G}_{m}\to G$, par descente à partir du cas déployé : là,
$T\simeq\mathbf{G}_{m}^{\Sigma}$ par les racines simples,
$\operatorname{Hom}(\mathbf{G}_{m},T)\simeq\mathbf{Z}^{\Sigma}$, et l'on prend
l'élément dont toutes les coordonnées valent $1$ — soit $\rho^{\vee}$. Son
centralisateur est $T$. Un homomorphisme est \emph{principal} s'il est associé
à un couple de Killing ; ce couple est alors unique ($i$ est régulier), la
condition est locale pour fppf, et les homomorphismes principaux
correspondent bijectivement aux couples de Killing.

\textbf{(i bis)} Si $i$ est principal, $\lambda\mapsto i(\lambda^{-1})$ l'est
aussi, associé au couple de Killing opposé. Le normalisateur de
$i(\mathbf{G}_{m})$ est l'image inverse $N_{1}(T)$, dans $N(T)$, du sous-groupe
d'ordre $2$ de $W$ engendré par l'élément $w_{0}$ qui envoie $B$ sur le Borel
opposé (on suppose $r\geqslant 1$). Un sous-tore de rang $1$ est
\emph{principal} s'il est localement, pour fppf, l'image d'un homomorphisme
principal ; ces sous-tores correspondent aux sections de
$\mathrm{Kill}/(\mathbf{Z}/2\mathbf{Z})$, le quotient étant pris par $w_{0}$ :
un tore maximal et une paire de Borels opposés le contenant, ces Borels
n'étant pas nécessairement définis séparément sur $S$.

\pagerange{44}{46}
\subsection*{(ii) Homomorphismes de Kostant, et la graduation de $\mathfrak{g}$}

Supposons $G$ simple.\footnote{« $G$ simple ? » est en marge : il le faut pour
que le nombre de Coxeter $h$ soit défini.} Un homomorphisme principal $i$
induit $\mu_{h}\to G$. Un homomorphisme $K\colon\mu_{h}\to G$ est \emph{de
Kostant} s'il est localement (fppf, ou étale fini, cela revient au même)
obtenu ainsi. Il est alors \emph{régulier} : son centralisateur connexe est un
tore maximal $T$. La $\mathbf{Z}/h\mathbf{Z}$-graduation de $\mathfrak{g}$
qu'il définit a pour composante de degré $0$ la sous-algèbre de Cartan
$\mathfrak{t}$, et pour composante de degré $j$, $1\leqslant j\leqslant h-1$,
\[
  \mathfrak{g}^{j}=\sum_{\operatorname{ht}(\alpha)=j}\mathfrak{g}_{\alpha}
  \;+\sum_{\operatorname{ht}(\alpha)=j-h}\mathfrak{g}_{\alpha},
\]
de rang le nombre de racines dont la hauteur est congrue à $j$ modulo
$h$.\footnote{Le feuillet 44 dit « racines positives dont l'ordre est $\equiv j$
(mod $h$), i.e.\ dont l'ordre est $j$ ou $j-h$ » ; les racines de hauteur
$j-h$ sont négatives, et la formule du feuillet 45, qui somme sur toutes les
racines, est la bonne. Elle y est écrite pour $1\leqslant j\leqslant h$, mais
le degré $h$ est le degré $0$.} En particulier $\mathfrak{g}^{1}$, somme des
espaces des racines simples et de celui de $-\psi$, est de rang $r+1$.

Le centralisateur $G^{K}=N(T)^{K}$ est l'image inverse du stabilisateur $W_{K}$
de $K$ dans $W$. Kostant dit que l'ordre de $W_{K}$ est celui du centre de
$\widetilde{G}$, et Grothendieck espère qu'il soit isomorphe à $P/R$ (poids
modulo racines) ou à son dual à valeurs dans $\mathbf{Q}/\mathbf{Z}$. C'est la
seconde forme qui est juste : $W_{K}$ est isomorphe à $P^{\vee}/Q^{\vee}$,
dual de $P/Q$.\footnote{Le feuillet écrit « le rang de $W_{K}$ » pour son
ordre. En termes actuels, $K$ correspond au point $\rho^{\vee}/h$, intérieur à
l'alcôve fondamentale où toutes les racines affines simples valent $1/h$ ; son
stabilisateur dans le groupe de Weyl affine étendu est le groupe $\Omega$ des
symétries du diagramme de Dynkin étendu, isomorphe à $P^{\vee}/Q^{\vee}$, et
s'injecte dans $W$ avec pour image $W_{K}$. La formulation est de l'édition.}

La donnée de $K\colon\mu_{h}\to T$ équivaut à celle de l'ensemble $\Sigma'$ des
racines $\alpha$ telles que $\mathfrak{g}_{\alpha}\subset\mathfrak{g}^{1}$, au
moins si $h>2$ ; dans le cas déployé, $\Sigma'=\Sigma\cup\{-\psi\}$, la base
des racines affines simples. Grothendieck vérifie, dans le cas déployé, que
$W_{K}$ est le stabilisateur de $\Sigma'$ dans $W$, c'est-à-dire l'ensemble des
$w$ qui envoient toutes les racines simples dans $\Sigma$, sauf au plus une,
envoyée sur $-\psi$.\footnote{Le feuillet dit « sauf une » ; l'élément neutre
envoie $\Sigma$ dans $\Sigma$, d'où « au plus une ». L'inclusion $W_{K}\subset
\operatorname{Stab}(\Sigma')$ est notée « a priori ». La phrase sur
$\Sigma'$, avec ses mots « $\mathbf{Q}$-simples » et « $\mathbf{Q}$-base »,
est en partie illisible.} Il reste à déterminer le normalisateur de $K$, et
l'action de $\operatorname{Norm}(K)/\operatorname{Cent}(K)$ sur $\mu_{h}$ :
est-elle transitive sur $\operatorname{Aut}(\mu_{h})=
(\mathbf{Z}/h\mathbf{Z})^{*}$ ? Le dossier ne le dit pas.

\pagerange{46}{46}
\subsection*{(iii) Le théorème de Kostant sur $\mathfrak{g}^{1}$}

Soit $K\colon\mu_{h}\to T$ de Kostant et
$\mathfrak{g}^{1}=\sum_{\alpha\in\Sigma'}\mathfrak{g}_{\alpha}$. En
caractéristique nulle (Kostant) :

\begin{enumerate}
\item[a)] une section $\sum x^{\alpha}$ de $\mathfrak{g}^{1}$, dont aucune
  composante ne s'annule, est semi-simple régulière ;
\item[b)] si l'une des composantes est nulle, elle est nilpotente ;
\item[c)] deux droites $D$, $D'$ du cas a) sont conjuguées par un $t\in T$,
  sans unicité : le stabilisateur dans $T$ d'une telle droite est $K(\mu_{h})$,
  et les droites du cas a) forment un torseur sous $T/K(\mu_{h})$.
\end{enumerate}

Les trois points se voient sur les poids : $T$ agit sur les $r+1$
coordonnées $x^{\alpha}$ par les caractères $\alpha\in\Sigma'$, liés par la
seule relation $-\psi+\sum_{i}c_{i}\alpha_{i}=0$, où les $c_{i}$ sont les
coefficients de $\psi$ ; les invariants de $\mathfrak{g}$ restreints à
$\mathfrak{g}^{1}$ sont donc des polynômes en
$x^{-\psi}\prod_{i}(x^{\alpha_{i}})^{c_{i}}$, qui s'annule dès qu'une
coordonnée s'annule, d'où b). Et $t$ stabilise une droite du cas a) si et
seulement si tous les $\alpha(t)$, $\alpha\in\Sigma'$, sont égaux à un même
$\lambda$, ce qui force $\lambda^{1+\sum c_{i}}=\lambda^{h}=1$ et
$t=K(\lambda)$.\footnote{Cette vérification est de l'édition ; elle établit
aussi l'énoncé « à prouver » du feuillet 30. La page ajoute, à l'encre brune et
en partie illisiblement, qu'une identification de $D$ à $D'$ dans la bonne
classe est induite par un unique $t\in T$. Dans c), « loc. » est de lecture
incertaine.}

\pagerange{47}{47}
\subsection*{(iv) Droites cycliques}

Une \emph{droite cyclique} est une droite $D$ de $\mathfrak{g}$ qui, localement
pour fppf, est propre de caractère $\chi_{1}$ sous un homomorphisme de Kostant
$\mu_{h}\to G$ (et du cas a)). Les droites cycliques sont conjuguées. Si $D$
est cyclique, son centralisateur connexe est un tore maximal $T'$, et son
normalisateur est un sous-groupe de $N(T')$ contenant $T'$, image inverse
d'un sous-groupe de $W'=N(T')/T'$ canoniquement isomorphe à $\mu_{h}$, dont un
générateur correspond à un élément de Coxeter.\footnote{La page assortit
l'énoncé sur le centralisateur d'une réserve, « (sous réserve que ce soit
unique) », de lecture incertaine. Les mots « fini » (biffé) et « can.\ isom.\
à » sont interlinéaires. Le reste du feuillet est blanc.}

C'est, relevé dans $G$ et sur une base, le dernier corollaire des feuillets 1
et 2 : le dossier se referme sur l'objet par lequel il s'était ouvert.

\end{document}
